% !TeX root = ../../Book1.tex
\section{可数性与乘积拓扑}
\label{b1:sec:A103}

可分性提供一列测试点，第二可数性提供一列测试开集。前者谈“点够不够多”，后者谈“邻域能否由可数种形状生成”。在度量空间中它们等价，在一般拓扑中却不能混用。中文教材中这两个名称可对照熊金城《点集拓扑讲义》第五版\cite[第5章]{Xiong2020Topology}；其积空间部分也适合作为本节的进一步阅读。

\subsection{从稠密点到可数基}
\label{b1:subsec:A10301}

\begin{definition}\label{b1:def:app-countability}
若 \(D\subset X\) 满足 \(\overline D=X\)，则称 \(D\) 在 \(X\) 中稠密；若存在可数稠密子集，则称 \(X\) \term[kefen]{可分}[separable]。若开集族 \(\mathcal G\) 满足每个开集均是其中若干成员的并，则称它为拓扑基。若有可数拓扑基，则称 \(X\) 满足\term[dier-keshuxing]{第二可数性}[second countability]。
\end{definition}

拓扑基的等价用法是：对 \(x\in U\)、\(U\) 开，可取 \(G\in\mathcal G\)，使 \(x\in G\subset U\)。稠密集则只要求与每个非空开集相交。两者之间少了对开集大小的控制，所以一般不能直接用稠密集的点替代一组基。

\begin{theorem}\label{b1:thm:app-separable-base}
第二可数空间可分，而且每个开覆盖都有可数子覆盖。对度量空间，可分性又推出第二可数性。第二可数空间的每个子空间仍第二可数；因此可分度量空间的每个子空间可分。
\end{theorem}
\begin{proof}
从可数基的每个非空成员中取一点，所得可数集与每个非空开集相交，故稠密。给定开覆盖，只考虑那些包含在某一覆盖成员中的基本开集，并为每个这样的基元选一个包含它的成员。这些选择至多可数，且覆盖全空间。

设 \((X,d)\) 有可数稠密集 \(D\)。以 \(D\) 中的点为中心、正有理数为半径的球组成可数族。给定 \(x\in U\)，先取 \(\varepsilon>0\)，使 \(B_d(x,\varepsilon)\subset U\)，再取 \(a\in D\) 满足 \(d(a,x)<\varepsilon/4\)，取有理数 \(r\) 满足 \(d(a,x)<r<\varepsilon/2\)。于是 \(x\in B_d(a,r)\subset U\)，这证明它是一组基。

若 \(\mathcal G\) 为 \(X\) 的可数基，则 \(\{A\cap G:G\in\mathcal G\}\) 是 \(A\) 的可数基。最后一句由此前两项应用于度量子空间得到。
\end{proof}

最后一句不能用 \(D\cap A\) 来证明：稠密集与子空间的交可能为空，例如 \(\mathbb Q\cap(\mathbb R\setminus\mathbb Q)\)。真正继承给子空间的是可数基，再由新基选择稠密点。

有限球网也提供可数点集。全有界空间对每个 \(n\) 取一个有限 \(1/n\)-网，把所有中心取并便得可数稠密集。因此紧度量空间可分；它的连续函数可常常通过可数点上的信息确定，但函数之间的一致估计仍需要连续性或等度连续性，不能仅凭“可数”二字省去。

\begin{proposition}\label{b1:prop:app-continuous-dense}
若 \(D\) 在 \(X\) 中稠密，\(f,g:X\rightarrow Y\) 连续，且 \(Y\) 为 Hausdorff 空间，则 \(f=g\) 只需在 \(D\) 上成立。若 \(X\) 可分，则连续像 \(f(X)\) 也可分。
\end{proposition}
\begin{proof}
若 \(f(x)\ne g(x)\)，用互不相交的开邻域分开这两个点，其反像的交给出 \(x\) 的开邻域，在其中 \(f\) 与 \(g\) 处处不同，与该邻域遇到 \(D\) 矛盾。第二句取可数稠密集 \(D\)；像空间中的非空相对开集有非空开反像，反像遇到 \(D\)，因此该开集遇到 \(f(D)\)。
\end{proof}

\subsection{有限坐标控制}
\label{b1:subsec:A10302}

\secref{b1:sec:0903}已经引入乘积拓扑并证明 Tychonoff 定理。这里补充其使用方法。设 \(X=\prod_{i\in I}X_i\)，坐标投影为 \(\pi_i\)。基本开集是
\[
 \bigcap_{i\in F}\pi_i^{-1}(U_i),\quad F\subset I\text{ 有限},\quad U_i\subset X_i\text{ 开}.
\]
不在 \(F\) 中的坐标完全不受限制。特别地，乘积拓扑是使全部坐标投影连续的最弱拓扑：任何这样的拓扑都必须包含上述有限交及其任意并。

\begin{proposition}\label{b1:prop:app-product-universal}
映射 \(f:Z\rightarrow\prod_iX_i\) 连续，当且仅当每个 \(\pi_i\circ f\) 连续。一个网在乘积拓扑中收敛，当且仅当它在每个坐标中收敛。
\end{proposition}
\begin{proof}
必要性由投影连续得到。充分性只需考察基本开集的反像：它是有限个坐标反像的交，因而开。对网的收敛，给定一个限制有限坐标的基本邻域，每个受限坐标都在某个指标以后满足要求。有向性保证这有限个指标有共同上界，于是从该上界起同时满足全部限制。
\end{proof}

若把“有限坐标”换成“所有坐标”，得到的更强拓扑一般不再有这一结论。例如 \(\mathbb R^{\mathbb N}\) 中的标准单位向量 \(e_n\) 逐坐标趋零，但在要求所有坐标同时属于 \((-1/2,1/2)\) 的盒状邻域中，没有一个 \(e_n\) 属于它。这个集合不是乘积拓扑的零点邻域。

\begin{theorem}[可数乘积的度量]\label{b1:thm:app-countable-product}
设 \((X_j,d_j)\)、\(j\ge1\)，是非空度量空间。下式给出其乘积拓扑的一个度量：
\[
 D(x,y)=\sum_{j=1}^{\infty}2^{-j}\min\{1,d_j(x_j,y_j)\}.
\]
若每个 \(X_j\) 完备，则 \(D\) 完备；若每个 \(X_j\) 可分，则乘积可分；若每个 \(X_j\) 紧，则乘积紧。
\end{theorem}
\begin{proof}
截断距离仍满足三角不等式，且上述级数绝对收敛，所以 \(D\) 是度量。给定限制有限坐标的基本邻域，对这些坐标取相应正距离阈值，再选足够小的 \(D\)-球，便落在其中。反之，给定 \(\varepsilon>0\)，先选 \(N\)，使 \(\sum_{j>N}2^{-j}<\varepsilon/2\)，再把前 \(N\) 个坐标限制在足够小的球中，使级数前段小于 \(\varepsilon/2\)。所得基本邻域包含于 \(D\)-球。故两种拓扑一致。

若 \((x^{(n)})\) 为 \(D\)-Cauchy 列，固定坐标 \(j\) 后，由其对应的一项小于总和，得到 \((x_j^{(n)})\) 是 \(d_j\)-Cauchy 列。逐坐标取极限 \(x_j\)，再用前段有限和与尾和估计，得 \(D(x^{(n)},x)\rightarrow0\)。

为证可分，对每个因子选可数稠密集 \(A_j\) 及基点 \(a_j\in A_j\)。考虑那些只有有限个坐标偏离 \(a_j\)，且每个坐标属于 \(A_j\) 的点。这个点集可数，并与每个非空基本开集相交。紧性可由\cref{b1:thm:tychonoff}直接得到；也可逐坐标抽收敛子列，再取对角子列，用上述尾和估计证明在 \(D\) 下收敛，最后应用度量紧性判据。
\end{proof}

这里可数个坐标使级数权重与对角子列都可用。一般不可数乘积的紧性仍成立，但不能从这一度量公式或可数次抽取推出，这正是 Tychonoff 定理在\chapref{b1:ch:09}中不可被简单对角法替代的原因。

\subsection{Borel 结构中的可数性}
\label{b1:subsec:A10303}

\begin{proposition}\label{b1:prop:app-product-borel}
若每个 \(X_j\) 第二可数，\(j\ge1\)，则
\[
 \mathcal B\left(\prod_{j=1}^{\infty}X_j\right)
 =\bigotimes_{j=1}^{\infty}\mathcal B(X_j),
\]
右侧指由各坐标投影的 Borel 反像生成的 \(\sigma\)-代数。
\end{proposition}
\begin{proof}
投影连续，故右侧包含于左侧。反向包含只需证明每个开集属于右侧。各因子选择可数基，由有限个坐标基元组成的基本开集只有可数多个，它们组成乘积的一组可数基。每个基本开集是有限个坐标 Borel 反像的交，每个开集是至多可数个这类集合的并，故属于右侧。
\end{proof}

证明中不是“开集是基本开集之并”就已足够，因为 \(\sigma\)-代数只对可数并封闭。可数基把任意并变成可数并，这是拓扑假设真正进入测度论的位置。\chapref{b1:ch:07}的乘积测度使用可测矩形，不能在一般空间中不经这一步就把乘积可测集与全部乘积 Borel 集等同。
